\documentclass[journal]{IEEEtran}
\IEEEoverridecommandlockouts
\usepackage{cite}
\usepackage{amsmath,amssymb,amsfonts}
\usepackage{algorithm}
\usepackage{algorithmic}
\usepackage{imakeidx}
\usepackage{booktabs}
\usepackage{graphicx}
\usepackage{multirow}
\usepackage{float}
\usepackage{subfigure}
\usepackage{bbding}
\usepackage{pifont}
\usepackage{textcomp}
\usepackage{xcolor}
\usepackage{hyperref}
\usepackage{tabularx}
\usepackage{comment}
\usepackage{array}
\usepackage{makecell}
\usepackage{placeins}

\usepackage[capitalise,noabbrev]{cleveref}
\Crefname{figure}{\text{Fig.}}{\text{Figs.}}
\Crefname{equation}{\text{Eq.}}{\text{Eqs.}}
\def\BibTeX{{\rm B\kern-.05em{\sc i\kern-.025em b}\kern-.08em
    T\kern-.1667em\lower.7ex\hbox{E}\kern-.125emX}}
\usepackage{geometry}
\geometry{left=0.75in,right=0.75in,top=0.75in,bottom=0.75in}

% ============================================================================
% REVISION NOTE (for the authors, not for submission -- remove before upload):
% This draft has been reframed for a journal submission (IEEE T-ITS).
% Key changes vs. the rejected conference version:
%  (1) The flow is now conditioned on a deterministic latent WORLD-MODEL
%      rollout instead of a raw time index -> temporally coherent risk
%      tensor (answers Reviewer 2's per-step-independence critique).
%  (2) Action-conditioned COUNTERFACTUAL risk field -> a capability PORA
%      (Wang et al., arXiv:2501.16480) structurally lacks (answers Reviewer 1
%      / AE novelty critique).
%  (3) Experiments restructured to compare against OTHER RISK QUANTIFICATION
%      METHODS (TTC / extended-TTC / DSF / EDRF / RiskNet / PORA-style
%      occupancy overlap), per the AE's "meaningful baseline" requirement.
%      Trajectory-prediction baselines (FloMo etc.) intentionally dropped
%      since prediction accuracy is not the contribution.
%  (4) Duplicate Vemula "Social Attention" reference removed.
%  TODO markers (\TODO{...}) flag every quantity that must be filled from a
%  real experimental run -- numbers have NOT been fabricated.
% ============================================================================
\newcommand{\TODO}[1]{{\color{red}\textbf{[TODO: #1]}}}

\begin{document}
\newgeometry{left=0.75in,right=0.75in,top=1in,bottom=0.75in}

\bstctlcite{BSTcontrol}

\title{World-Model-Conditioned Normalizing Flows for Spatial-Temporal Risk Assessment of Autonomous Vehicles}
\author{
}

\maketitle

\begin{abstract}
Safety assessment of autonomous vehicles needs risk representations that are predictive, temporally coherent, interpretable, and able to answer \emph{what-if} questions about the ego's own decisions. Existing approaches often fall short in these and more importantly cannot provide the risk evolution over time. We propose an ego-centric risk field built on \emph{world-model-conditioned} normalizing flows: a deterministic latent state rolled over the horizon conditions an invertible flow, giving an exact, grid-evaluable occupancy density without sampling. The field couples two such flows, a single-target model for the ego's occupancy and a scene-level \emph{autoregressive joint} for the surrounding agents conditioned on the ego, both conditioned on a bird's-eye \emph{map} of the road. Per-cell risk is the ego occupancy times an oriented bounding-box meeting probability times a physically grounded kinetic-energy severity, with relative velocities from the centroid motion of the predicted density. On InD intersections and the congested-highway AD4CHE dataset, the field detects conflicts competitively-to-better than kinematic and occupancy-overlap baselines, with the largest margin in congested traffic and supports a \emph{physically grounded counterfactual} that fixed-plan methods structurally lack. Our method provides anticiapted risk representations, enabling predictive safety beyond state-of-the-art.
\end{abstract}

\begin{IEEEkeywords}
Autonomous vehicles, safety assessment, risk field, world model, counterfactual reasoning, conditional normalizing flows, probabilistic occupancy
\end{IEEEkeywords}

\section{Introduction}\label{section:introduction}
Safety assessment and validation remain central challenges for autonomous vehicles (AVs), especially in dense, interactive traffic where multiple plausible futures may unfold. Converting multi-agent uncertainty into an actionable safety representation is still non-trivial: practitioners need a representation that is (i) spatially continuous, (ii) predictive over a time horizon, (iii) interpretable enough to support engineering decisions, (iv) \emph{temporally coherent}, so that the risk shown at successive horizon steps corresponds to a consistent future rather than a sequence of unrelated snapshots, and (v) \emph{counterfactual}, so that the effect of the ego vehicle's own candidate decisions on the risk landscape can be evaluated before acting.

A common line of work represents risk as continuous fields. Inspired by artificial potential fields \cite{xie2022distributed}, traffic-oriented safety fields map hazards into spatial potentials or ``field strengths'' that reflect driver--vehicle--road interactions, with extensions that aim to reproduce human-like risk sensitivity. These field-based approaches are intuitive and easy to visualize, but their functional forms and parameters are often handcrafted, which limits their ability to cover complex, multi-modal interactions. Classical surrogate safety measures (SSMs) such as time-to-collision (TTC) and its variants provide efficient and interpretable conflict indicators and underpin system-level safety analysis, but rely on short-horizon kinematic assumptions and do not yield a probabilistic description of where risk may materialize over space and time.

Learning-based motion prediction models can produce multi-modal forecasts by modeling interactions. For safety assessment, a recurring difficulty is that many predictors output sampled trajectories, anchors, or discrete occupancies, which require additional post-processing to obtain a spatially continuous risk representation. Closest to our setting, recent spatio-temporal probabilistic occupancy-heatmap methods, in particular the Probabilistic Occupancy Risk Assessment (PORA)~\cite{wang2025dynamic}, generate a stack of bird's-eye-view occupancy heatmaps over future steps and adjust collision risk along a \emph{given} AV plan using occupancy changes. While powerful, these methods share three limitations that motivate our work: their occupancy is a learned discretized map rather than an exact normalized density; the per-step maps are produced largely independently, so temporal evolution is not guaranteed to be coherent; and conditioning on a single fixed AV plan does not support reasoning about alternative ego actions.

Normalizing flows offer an exact, grid-evaluable conditional density through an invertible transformation, removing the sampling/stitching post-processing. However, a flow conditioned only on a context vector and a raw time index still produces independent per-step marginals. We close this gap by borrowing the central idea of \emph{world models} : a learned latent state that is rolled forward through time and, crucially, conditioned on actions. We use a deterministic latent world model as the temporal backbone that conditions the flow, so successive density maps share an evolving state (temporal coherence), and we drive the rollout with the ego action sequence so the resulting field is counterfactual.

Concretely, the field is ego-centric: at each future step and grid location, risk is the ego's own occupancy likelihood times the probability that the ego's footprint meets a surrounding agent there times a physically grounded kinetic-energy-loss severity term, aggregated over agents into a horizon-indexed risk tensor. This couples two world-model-conditioned flows, a single-target model for the ego's occupancy and a scene-level autoregressive joint for the surrounding agents' occupancies \emph{conditioned on the ego}, both conditioned on a bird's-eye map of the road geometry.

Our main contributions are:
\begin{itemize}
\item We propose an invertible flow conditioned on a deterministic latent world-model rollout yields an exact, grid-evaluable occupancy density without sampling or trajectory stitching.
\item We couple a single-target flow for the ego's occupancy with a scene-level autoregressive joint for the surrounding agents \emph{conditioned on the ego}. Per-cell risk multiplies the ego occupancy, an oriented-box meeting probability, and a kinetic-energy severity.
\item On InD intersections and the congested-highway AD4CHE dataset, the field is competitive-to-best at conflict detection, with the largest margin on congested traffic where kinematic baselines fail.
\end{itemize}

The remainder of the paper is organized as follows. \Cref{section:related_work} reviews related work and positions our method against occupancy-heatmap risk assessment. \Cref{section:Methodology} presents the world-model-conditioned flow, the two-model (ego marginal and scene-level joint) field with map conditioning, and the physically grounded counterfactual risk decomposition. \Cref{sec:results} reports the result on the proposed approach, followed by conclusions and future works in \Cref{section:conclusion_future_work}.

\section{Related Work}\label{section:related_work}
This section reviews (i) field-based risk representations, (ii) classical surrogate safety measures, (iii) learning-based and occupancy-heatmap risk methods, and (iv) world models, and states the gap our method targets.

\subsection{Field-Based Methods}\label{sec:rw_field_based}
Field representations originate from artificial potential fields in robotics, where obstacles map to continuous potentials for real-time planning. In traffic safety, the Driving Safety Field (DSF) and related studies model driver--vehicle--road interactions and potential hazards as continuous safety/risk fields, with extensions for human-like risk-related behavior. These methods are spatially continuous and physically intuitive, but the field form and parameters are typically handcrafted, which makes it hard to cover complex, multi-modal interactions without extra assumptions, and they are not derived from a calibrated predictive distribution.

\subsection{Classical Risk Assessment}\label{sec:rw_classical}
Kinematic conflict indicators such as TTC and its variants are widely used for safety evaluation and warning. Together with measures like headway, they provide simple and interpretable surrogate signals of conflict severity, and are often used within the Surrogate Safety Measures (SSMs) framework for system-level analysis. However, they are usually computed under short-horizon assumptions (e.g., constant velocity/acceleration), have limited predictive power for long-horizon, interactive, multi-modal futures, and do not provide a probabilistic description of risk over space and time.

\subsection{Learning-Based and Occupancy-Heatmap Risk Methods}\label{sec:rw_multimodal}
Deep learning has enabled multi-modal trajectory prediction by modeling interactions and generating multiple plausible futures. Many methods output a fixed set of discrete trajectories or anchors, and probability is implicit rather than a normalized density; converting these into a spatially continuous risk field requires additional post-processing. A growing body of work instead predicts probabilistic occupancy directly: occupancy flow fields, enhanced driving risk fields from multimodal prediction (EDRF), uncertainty-aware probabilistic SOTIF quantification, interaction-aware risk forecasting (RiskNet), and, most relevant here, the spatio-temporal probabilistic occupancy heatmaps of PORA, which evaluate collision risk along a given AV plan and adjust it through occupancy changes and a Cox-process correction.

Our method shares the high-level goal of a spatio-temporal probabilistic risk representation but differs in three concrete ways. First, our occupancy term is an \emph{exact, normalized} conditional density evaluated analytically on the grid through an invertible flow, rather than a learned discretized heatmap; this removes binning/sampling noise and yields a mathematically transparent density. Second, our severity term is a \emph{physical} kinetic-energy-loss surrogate, in contrast to occupancy-overlap or statistical adjustments. Third, and most importantly, our per-step densities are conditioned on a \emph{shared latent world-model rollout} and on the \emph{ego action sequence}, which makes the risk tensor temporally coherent and \emph{counterfactual}; PORA and related occupancy methods condition on a single fixed AV plan and produce per-step maps that are not guaranteed to form a coherent future. \Cref{tab:positioning} summarizes this positioning.

\begin{table}[h]
\caption{Positioning against representative risk-quantification methods.}
\label{tab:positioning}
\centering
\footnotesize
\begin{tabularx}{\columnwidth}{l c c c c}
\toprule
\textbf{Property} & \textbf{SSM} & \textbf{DSF} & \textbf{PORA} & \textbf{Ours} \\
 & \cite{7,10} & \cite{4,6} & \cite{wang2025dynamic} & \\
\midrule
Spatially continuous        & \ding{55} & \ding{51} & \ding{51} & \ding{51} \\
Probabilistic               & \ding{55} & \ding{55} & \ding{51} & \ding{51} \\
Exact normalized density    & \ding{55} & \ding{55} & \ding{55} & \ding{51} \\
Temporally coherent rollout & \ding{55} & \ding{55} & partial   & \ding{51} \\
Action-conditioned (counterfactual) & \ding{55} & \ding{55} & \ding{55} & \ding{51} \\
Physical severity model     & partial   & partial   & \ding{55} & \ding{51} \\
\bottomrule
\end{tabularx}
\end{table}

\subsection{World Models}\label{sec:rw_worldmodel}
World models learn a latent state that is rolled forward through time and conditioned on actions, enabling planning and imagination in latent space. They have been adapted to urban driving for model-based imitation and forecasting. We do not use a world model to generate pixels or to act; we use a lightweight deterministic latent rollout purely as the temporal, action-conditioned backbone that conditions an exact-density flow. This is, to our knowledge, a new use of the world-model idea for constructing an interpretable, counterfactual risk field.


\section{Methodology} \label{section:Methodology}
The key idea is to decompose traffic risk into (i) a collision likelihood (density) map on a grid and (ii) the physical cost of a collision, where each density is produced by an invertible flow conditioned on a deterministic latent \emph{world-model} state rolled forward over the horizon, on a \emph{map} encoding of the road geometry, and on the ego action sequence. The risk field is \emph{ego-centric}: it is the risk the ego itself would face, so it combines \emph{two} world-model-conditioned flows that play complementary roles. A \emph{single-target marginal} model supplies the ego's own occupancy $P_{\mathrm{ego}}$ (and, by index rotation, any agent's self-motion); a \emph{scene-level autoregressive joint} model supplies the surrounding agents' occupancies \emph{conditioned on the ego}, $p(\mathbf{Y}_{1:N}\mid\text{scene},\mathbf{a}_{\mathrm{ego}})=\prod_i p(\mathbf{Y}_i\mid\mathbf{Y}_{<i},\text{scene},\mathbf{a}_{\mathrm{ego}})$. This yields a temporally coherent and counterfactual 3D risk tensor for a given history.

An overview is shown in \Cref{fig:overall_architecture}. Dashed lines denote functionalities inactive during risk assessment (e.g., sample generation).

\begin{figure*}[!t]
\centering
\includegraphics[width=\textwidth]{figures/1.png}
\caption{Overall architecture: an encoder summarizes multi-agent history into a latent state, a deterministic world model rolls the state forward conditioned on the ego action sequence, and a conditional normalizing flow turns each rolled-out state into an exact occupancy density that is combined with a kinetic-energy-loss severity term to form the spatial-temporal (and counterfactual) risk tensor.}
\label{fig:overall_architecture}
\end{figure*}

\subsection{Problem Setup and Notation}\label{sec:problem_setup}
We consider a multi-agent scene with up to $N$ vehicles indexed by $i\in\{0,\dots,N-1\}$, where index $0$ is the \emph{ego}. The history length is $T$ and the prediction horizon is $K$ steps.

For vehicle $i$, the 2D position at historical step $\tau\in\{1,\dots,T\}$ is $\mathbf{p}^{(i)}_{\tau}\in\mathbb{R}^2$, with per-step kinematic features $\mathbf{f}^{(i)}_{\tau}\in\mathbb{R}^{d_f}$ --- here $d_f{=}5$: heading, longitudinal and lateral velocity, and longitudinal and lateral acceleration --- and a time channel $t_{\tau}$ (each vehicle additionally carries a categorical type used for FiLM modulation, \Cref{sec:context_encoding}):
\begin{equation}
\mathbf{x}^{(i)}_{\tau}=\big[\mathbf{p}^{(i)}_{\tau},\;\mathbf{f}^{(i)}_{\tau},\;t_{\tau}\big].
\label{eq:input_def}
\end{equation}
The full history is $\mathbf{X}=\{\mathbf{x}^{(i)}_{1:T}\}_{i=0}^{N-1}$. We discretize the region of interest into an $S\times S$ grid of cell centers $\mathcal{G}=\{\mathbf{g}_{m,n}\}_{m,n=1,\dots,S}$ with spatial resolution $\Delta$ and sampling interval $\Delta t$, and we attach to each location a static bird's-eye \emph{map raster} $\mathbf{M}\in\mathbb{R}^{3\times S_m\times S_m}$ of the road geometry (\Cref{sec:map_conditioning}). We introduce an \emph{ego action sequence} $\mathbf{a}_{1:K}$, with $\mathbf{a}_k\in\mathbb{R}^{d_a}$ a compact control/intent descriptor at step $k$ (e.g., longitudinal acceleration and yaw rate, or a waypoint displacement).

\paragraph{Two predictive objects.} The ego-centric risk field requires two distinct quantities, each a world-model-conditioned flow density evaluated exactly on $\mathcal{G}$:
\begin{enumerate}
\item the ego's own future occupancy $P_{\mathrm{ego},k}(\mathbf{g})$, from a \emph{single-target marginal} model $p_{\theta}(\mathbf{y}_k\mid\mathbf{X},\mathbf{M})$ in which the predicted vehicle is placed at index $0$; by cyclically permuting the ordering so any vehicle $j$ becomes index $0$ (index-rotation inference, shared parameters) the same model supplies each agent's \emph{self-motion} marginal;
\item the surrounding agents' future occupancies \emph{conditioned on the ego's planned action}, from a \emph{scene-level autoregressive joint} model $p_{\psi}\!\big(\mathbf{Y}_{1:N}\mid\mathbf{X},\mathbf{M},\mathbf{a}_{1:K}\big)=\prod_i p_{\psi}(\mathbf{Y}_i\mid\mathbf{Y}_{<i},\mathbf{X},\mathbf{M},\mathbf{a}_{1:K})$ (\Cref{sec:scene_joint}), in which the ego is not predicted but supplies the conditioning action.
\end{enumerate}
The ego marginal $\theta$ is needed because the joint model leaves the ego unpredicted (its future is the imposed conditioning); the joint $\psi$ is needed because the risk the ego faces is shaped by where the \emph{others} go \emph{given the ego's plan}, not by their context-free marginals.

\subsection{Risk Decomposition}\label{sec:risk_decomposition}
The ego-centric risk at future step $k$ and grid point $\mathbf{g}\in\mathcal{G}$ is the chance the ego \emph{at} $\mathbf{g}$ meets a surrounding agent there, times the energy released if it does:
\begin{equation}
\mathrm{Risk}_{k}(\mathbf{g}\mid\mathbf{a}_{1:k})=P_{\mathrm{ego},k}(\mathbf{g})\cdot M^{(j)}_{k}(\mathbf{g}\mid\mathbf{a}_{1:k})\cdot C^{(j)}_{k}(\mathbf{g}),
\label{eq:risk_decomp}
\end{equation}
where $P_{\mathrm{ego},k}(\mathbf{g})$ is the ego's own occupancy density at $\mathbf{g}$ (\Cref{sec:collision_probability}), $M^{(j)}_{k}$ is the probability that an ego footprint at $\mathbf{g}$ overlaps agent $j$'s footprint under $j$'s ego-conditioned occupancy (\Cref{sec:box_overlap}), and $C^{(j)}_{k}$ is the collision energy if they meet (\Cref{sec:collision_cost}). Conditioning the surrounding agents on $\mathbf{a}_{1:k}$ makes the field \emph{counterfactual}: it is the risk the ego would face if it executed action sequence $\mathbf{a}_{1:k}$ (\Cref{sec:counterfactual}). Aggregating over surrounding vehicles $\mathcal{J}$:
\begin{equation}
\mathrm{Risk}^{\mathrm{all}}_{k}(\mathbf{g}\mid\mathbf{a}_{1:k})=P_{\mathrm{ego},k}(\mathbf{g})\sum_{j\in\mathcal{J}} M^{(j)}_{k}(\mathbf{g}\mid\mathbf{a}_{1:k})\,C^{(j)}_{k}(\mathbf{g}).
\label{eq:risk_sum}
\end{equation}
Unlike a product of two independent marginals $P_{\mathrm{ego}}\cdot P^{(j)}$, the surrounding density here is drawn from the joint conditioned on the ego (and on the other agents through the chain rule), so the field expresses the ego's risk \emph{given its own plan}.

\subsection{Latent World Model}\label{sec:world_model}
We replace the conference-version conditioning on a raw time index $k$ with a deterministic latent world-model rollout, which is the source of temporal coherence and counterfactual capability.

\subsubsection{Multi-Agent Context Encoding}\label{sec:context_encoding}
Each vehicle history is encoded by a temporal encoder (GRU):
\begin{equation}
\mathbf{h}^{(i)}=\mathrm{GRU}_{\theta}\big(\mathbf{x}^{(i)}_{1:T}\big)\in\mathbb{R}^{d_h}.
\label{eq:gru_encode}
\end{equation}
To capture heterogeneous vehicle characteristics, $\mathbf{h}^{(i)}$ is modulated by a FiLM-style vehicle-type embedding (per-channel scale and shift). The encoder has two interaction heads that share the per-agent histories but differ in how agents attend to one another, matching the two predictive objects of \Cref{sec:problem_setup}:
\begin{itemize}
\item \textbf{Single-target (marginal) head.} Multi-head cross-attention with the index-$0$ target as Query and the others as Key/Value, with a key-padding mask for absent neighbours, yields a permutation-invariant interaction vector $\mathbf{a}^{\mathrm{int}}$; a small MLP projects $[\mathbf{h}^{(0)},\mathbf{a}^{\mathrm{int}}]$ into a context vector $\mathbf{c}\in\mathbb{R}^{d_c}$ for the predicted agent.
\item \textbf{Scene-level (joint) head.} Symmetric multi-head self-attention over all valid agents (masked), with a learned \emph{ego-identity} vector added to index~$0$ so the ego remains distinguishable under the otherwise permutation-equivariant attention; this yields a per-agent embedding $\mathbf{c}^{(i)}\in\mathbb{R}^{d_c}$ for every agent, used by the autoregressive joint of \Cref{sec:scene_joint}.
\end{itemize}

\subsubsection{Map Conditioning}\label{sec:map_conditioning}
Public trajectory datasets such as InD provide a calibrated drone orthophoto but no vectorized lanes. We therefore resample, once per location, the orthophoto into the same normalized coordinate box the agents live in, giving a per-location bird's-eye raster $\mathbf{M}\in\mathbb{R}^{3\times S_m\times S_m}$ of the drivable area. A small convolutional encoder maps it to a map embedding
\begin{equation}
\mathbf{m}=E_{\mathrm{map}}(\mathbf{M})\in\mathbb{R}^{d_c},
\label{eq:map_embed}
\end{equation}
which is added to the agent embeddings ($\mathbf{c}\!\leftarrow\!\mathbf{c}+\mathbf{m}$ and $\mathbf{c}^{(i)}\!\leftarrow\!\mathbf{c}^{(i)}+\mathbf{m}$) before the world-model rollout, so the predicted densities follow the road geometry. The map is static per location, so it is precomputed and looked up by location at run time. As \Cref{sec:results_ablation} shows, this conditioning is load-bearing: removing it at inference sharply degrades held-out likelihood.

\subsubsection{Deterministic Latent Rollout}\label{sec:latent_rollout}
The (map-augmented) context initializes the latent state, $\mathbf{s}_0=\psi_{\phi}(\mathbf{c})$, which is rolled forward by a deterministic transition $g_{\phi}$ (a GRU/MLP cell) conditioned on a per-step control input $\mathbf{u}_k$:
\begin{equation}
\mathbf{s}_k=g_{\phi}\big(\mathbf{s}_{k-1},\,\mathbf{u}_k\big),\qquad k=1,\dots,K.
\label{eq:wm_transition}
\end{equation}
For the single-target marginal model, $\mathbf{u}_k=\mathbf{a}_k$ when the index-$0$ vehicle is the ego and a learned autonomous input $\mathbf{u}_k=\mathbf{0}$ otherwise (the transition then models typical maneuver dynamics conditioned on $\mathbf{c}$). For the scene-level joint model, the per-agent embeddings are pooled (masked mean, with the ego's embedding retained by a residual) into a single \emph{scene} state $\mathbf{s}_0^{\mathrm{sc}}$ that is rolled by the ego action, $\mathbf{s}_k^{\mathrm{sc}}=g_{\phi}(\mathbf{s}_{k-1}^{\mathrm{sc}},\mathbf{a}_k)$, so that changing the ego plan propagates to every agent's conditioning. Because $\mathbf{s}_k$ is produced by a recurrent transition from $\mathbf{s}_{k-1}$, consecutive horizon steps share an evolving state rather than being independent functions of $k$; this is precisely what makes the resulting density stack temporally coherent.

\subsection{Ego Occupancy via World-Model-Conditioned Flow}\label{sec:collision_probability}
We first describe the single-target marginal flow that supplies the ego occupancy $P_{\mathrm{ego}}$ (index $0$ is the ego; by index rotation the same flow gives any agent's self-motion). We define
\begin{equation}
P_{\mathrm{ego},k}(\mathbf{g})\equiv p_{\theta}\big(\mathbf{y}_k\mid\mathbf{s}_k\big)\big|_{\mathbf{y}_k=\mathbf{g}},
\label{eq:pk_def}
\end{equation}
where $\mathbf{s}_k$ is the world-model state from \Cref{eq:wm_transition}. The same flow architecture, conditioned instead on the scene-level autoregressive state of \Cref{sec:scene_joint}, supplies the ego-conditioned occupancies of the surrounding agents.

\subsubsection{Conditional Normalizing Flow Objective}\label{sec:flow_objective}
For each step $k$ we model $\mathbf{y}_k$ via an invertible map to a base Gaussian, conditioned on the rolled-out state:
\begin{equation}
\mathbf{z}_k=f_{\theta}\big(\mathbf{y}_k;\,\mathbf{s}_k\big),\qquad \mathbf{z}_k\sim\mathcal{N}(\mathbf{0},\mathbf{I}).
\label{eq:flow_forward}
\end{equation}
The conditional log-likelihood follows from change-of-variables:
\begin{equation}
\log p_{\theta}(\mathbf{y}_k\mid\mathbf{s}_k)=\log p(\mathbf{z}_k)+\log\left|\det\frac{\partial f_{\theta}}{\partial \mathbf{y}_k}\right|.
\label{eq:change_of_variables}
\end{equation}
The world model and flow are trained jointly by minimizing the horizon negative log-likelihood along the ground-truth-action rollout:
\begin{equation}
\mathcal{L}_{\mathrm{NLL}}(\theta,\phi)=-\sum_{k=1}^{K}\,\mathbb{E}\big[\log p_{\theta}\big(\mathbf{y}_k\mid \mathbf{s}_k(\mathbf{c},\mathbf{u}_{1:k})\big)\big].
\label{eq:nll}
\end{equation}
We use an invertible flow (affine coupling layers with invertible normalization). Note that $\prod_{k} p_{\theta}(\mathbf{y}_k\mid\mathbf{s}_k)$ with $\mathbf{s}_k$ a coherent latent rollout is a state-space factorization of the joint horizon distribution, not a product of independent marginals; the conference version is recovered by setting $\mathbf{s}_k\!=\![\mathbf{c},k]$ (no recurrence), which we use as the ablation in \Cref{sec:results_ablation}.

\subsubsection{Discretization and Footprint}\label{sec:grid_discretization}
At each step $k$ we evaluate the conditional density on the grid centers:
\begin{equation}
\rho_k(m,n)=p_{\theta}\big(\mathbf{y}_k\mid\mathbf{s}_k\big)\big|_{\mathbf{y}_k=\mathbf{g}_{m,n}}.
\label{eq:rho_grid}
\end{equation}
Per-cell mass is $\pi_k(m,n)\approx\rho_k(m,n)\,\Delta^2$; the constant $\Delta^2$ does not affect relative risk and is omitted. A vehicle occupies a finite footprint rather than a point; the simplest account convolves the point density with a fixed kernel $\mathbf{K}$,
\begin{equation}
P_k=\mathbf{K}\ast\rho_k,
\label{eq:kernel_dilation}
\end{equation}
which we generalize to an \emph{oriented} bounding-box overlap in \Cref{sec:box_overlap}. We write the (normalized) occupancy of agent $j$ as $P^{(j)}_{k}(\mathbf{g})$; for the ego this is $P_{\mathrm{ego},k}$, and for surrounding agents it is the ego-conditioned occupancy of \Cref{sec:scene_joint}.

\subsection{Joint Ego-Conditioned Occupancy of Surrounding Agents}\label{sec:scene_joint}
The risk the ego faces depends on where the \emph{other} agents go given the ego's plan. We therefore model the surrounding agents with a scene-level autoregressive joint, factorized by the chain rule under a fixed ordering $\pi$ (the ego at slot $0$, the remaining agents sorted by distance to the ego at the last history step, nearest first):
\begin{equation}
\begin{split}
p_{\psi}\big(&\mathbf{Y}_{\pi(1)},\dots,\mathbf{Y}_{\pi(N-1)}\mid\mathbf{X},\mathbf{M},\mathbf{a}_{1:K}\big)\\
&=\prod_{i\ge 1} p_{\psi}\big(\mathbf{Y}_{\pi(i)}\mid\mathbf{Y}_{\pi(<i)},\,\mathbf{s}^{\mathrm{sc}}_{1:K}\big).
\end{split}
\label{eq:scene_joint}
\end{equation}
Each conditional is the \emph{same} 2D flow as \Cref{eq:flow_forward}, evaluated exactly on the grid, but its per-step conditioning is produced by an autoregressive decoder
\begin{equation}
\mathbf{c}^{(i)}_k = \mathbf{c}^{(\pi(i))} + r_{\psi}\big(\mathbf{s}^{\mathrm{sc}}_k,\,\mathbf{c}^{(\pi(i))},\,\mathbf{h}^{\mathrm{prev}}_{k}(\mathbf{Y}_{\pi(<i)})\big),
\label{eq:ar_cond}
\end{equation}
where $\mathbf{s}^{\mathrm{sc}}_k$ is the ego-action-driven scene state (\Cref{sec:latent_rollout}) and $\mathbf{h}^{\mathrm{prev}}_k$ is a causal summary of the already-decoded agents' futures (a per-step cumulative pass, so $\mathbf{c}^{(i)}_k$ depends only on $\mathbf{Y}_{\pi(<i)}$). The residual $r_{\psi}$ has a zero-initialized last layer, so at initialization $\mathbf{c}^{(i)}_k=\mathbf{c}^{(\pi(i))}$ and the joint degrades to a per-agent baseline it can only refine. Training minimizes the masked per-agent horizon NLL over the surrounding agents (the ego is excluded, since its future is the conditioning), with teacher-forced $\mathbf{Y}_{\pi(<i)}$.

\paragraph{Inference (grid-exact, MAP rollout).} The first surrounding agent in $\pi$ conditions only on the ego, so its occupancy $P^{(\pi(1))}_k(\mathbf{g})$ is evaluated \emph{exactly} on the grid. For subsequent agents we fill $\mathbf{Y}_{\pi(<i)}$ with the per-frame MAP (grid argmax) of the already-decoded agents --- starting from the ego's MAP path, whose action proxy also forms $\mathbf{a}_{1:K}$ --- and evaluate $P^{(\pi(i))}_k$ on the grid given those MAPs. This preserves the exact, grid-evaluable density per conditional while keeping the chain tractable.

\subsection{Meeting Probability via Oriented Bounding-Box Overlap}\label{sec:box_overlap}
A collision is the overlap of two finite, oriented vehicle footprints, not the coincidence of two points. We model each vehicle as an oriented rectangle of type-dependent length and width, oriented by its per-cell heading (\Cref{sec:latent_velocity_field}; where an agent is near-stationary we fall back to its last observed yaw). The probability that an ego footprint placed at $\mathbf{g}$ overlaps agent $j$ is
\begin{equation}
M^{(j)}_{k}(\mathbf{g})=\sum_{\mathbf{g}'\in\mathcal{G}} P^{(j)}_{k}(\mathbf{g}')\,\mathbb{1}\!\big[\mathcal{B}_{\mathrm{ego}}(\mathbf{g})\cap\mathcal{B}_{j}(\mathbf{g}')\neq\varnothing\big],
\label{eq:box_overlap}
\end{equation}
where $\mathcal{B}$ denotes the oriented box and the intersection is decided by the separating-axis test (SAT) over the four box-edge normals. Implemented as a small set of per-orientation convolution kernels, \Cref{eq:box_overlap} reduces to \Cref{eq:kernel_dilation} when the kernel is the (isotropic) footprint. The oriented test removes false positives that an isotropic radius produces --- e.g.\ a car passing alongside a parked truck whose long axis never crosses the ego path.

\subsection{Collision Cost via Kinetic Energy Loss}\label{sec:collision_cost}
\subsubsection{Cost Definition}\label{sec:cost_definition}
We introduce a user-specified ego vehicle and compute cost ego-centrically; here $(0)$ denotes the ego. For surrounding vehicle $j\in\mathcal{J}$ and step $k$ we use a kinetic-energy-loss surrogate:
\begin{equation}
C^{(j)}_{k}=\tfrac{1}{2}\,m^{(j)}_{\mathrm{r}}\,\big\|\mathbf{v}^{(0)}_{k}-\mathbf{v}^{(j)}_{k}\big\|_2^2,
\label{eq:collision_cost}
\end{equation}
with reduced mass
\begin{equation}
m^{(j)}_{\mathrm{r}}=\frac{m^{(0)}m^{(j)}}{m^{(0)}+m^{(j)}}.
\label{eq:reduced_mass}
\end{equation}
Since public datasets (e.g., InD) lack masses, we assign a constant mass per vehicle category. The relative velocity is read from the predicted density's centroid motion (\Cref{sec:latent_velocity_field}), giving the severity $C^{(j)}_{k}$.

\subsubsection{Relative Velocity from the Density Centroid}\label{sec:latent_velocity_field}
The kinetic-energy cost needs the relative velocity of the two vehicles, which a static grid density does not carry; its \emph{evolution} does. The first moment (centroid) of the predicted occupancy traces each agent's bulk motion over the horizon: for agent $j$ we take the per-step occupancy centroid and its Savitzky--Golay-smoothed time derivative as the agent's velocity,
\begin{equation}
\hat{\mathbf{p}}^{(j)}_{k}=\!\!\sum_{\mathbf{g}\in\mathcal{G}}\!\mathbf{g}\,P^{(j)}_{k}(\mathbf{g}),\qquad
\mathbf{v}^{(j)}_{k}=\mathrm{SG}\big[\hat{\mathbf{p}}^{(j)}_{1:K}\big]_k\big/\Delta t,
\label{eq:centroid_topm}
\end{equation}
and obtain the ego velocity $\mathbf{v}^{(0)}_k$ identically from $P_{\mathrm{ego}}$. This centroid trajectory is temporally well-defined across steps, unlike the conference version's latent-code-preservation assumption, which is ill-defined for a flow trained on per-step marginals (the base code is not identified across steps). A per-cell alternative --- dense optical-flow (Horn--Schunck) transport between consecutive maps --- is location-dependent but over-reads at the spreading edge of the predicted density, where a small tail of front cells saturates the physical speed cap; since severity scales as $\lVert\mathbf{v}\rVert^2$, the energy peak then latches onto those outliers, so we use the robust bulk centroid instead. The severity is the relative kinetic energy,
\begin{equation}
C^{(j)}_{k}=\tfrac{1}{2}\,m^{(j)}_{\mathrm{r}}\,\big\|\mathbf{v}^{(0)}_{k}-\mathbf{v}^{(j)}_{k}\big\|_2^2,
\label{eq:collision_cost_field}
\end{equation}
a per-step relative kinetic energy applied over the meeting region. The heading used to orient the boxes in \Cref{eq:box_overlap} is the direction of $\mathbf{v}^{(j)}_k$ where the agent moves, and the last observed yaw where it is near-stationary.

\subsection{Counterfactual Risk Field}\label{sec:counterfactual}
Because the scene-level rollout (\Cref{sec:latent_rollout}) is driven by the ego action sequence, the entire risk tensor is a function of a candidate ego maneuver $\mathcal{A}=\mathbf{a}_{1:K}$ through \emph{two} channels: the surrounding agents' ego-conditioned occupancy $P^{(j)}_k(\cdot\mid\mathcal{A})$ (\Cref{eq:scene_joint}) and the kinetic-energy severity $C^{(j)}_k$ (\Cref{eq:collision_cost}), which depends on the ego's velocity profile under $\mathcal{A}$. The severity channel is causal by construction (braking lowers collision energy by mechanics); the density channel is a \emph{learned} reaction, and we report in \Cref{sec:results_negative} that on observational data it is weak --- the realized counterfactual signal is dominated by the physical severity term. Given a finite set of candidate maneuvers $\{\mathcal{A}^{(1)},\dots,\mathcal{A}^{(L)}\}$ (e.g., keep-lane, brake, yield, turn), we evaluate
\begin{equation}
J(\mathcal{A})=\sum_{k=1}^{K}\gamma^{k}\!\!\sum_{\mathbf{g}\in\mathcal{G}}\!\mathrm{Risk}^{\mathrm{all}}_{k}\big(\mathbf{g}\mid\mathcal{A}\big),
\label{eq:counterfactual_objective}
\end{equation}
with discount $\gamma\in(0,1]$, and report the per-maneuver risk landscape (and, optionally, $\mathcal{A}^\star=\arg\min_{\mathcal{A}} J(\mathcal{A})$ as a risk-aware action recommendation). Fixed-plan occupancy methods such as PORA~\cite{wang2025dynamic} evaluate only the single supplied plan and cannot produce this comparison.

\subsection{Risk Field Construction}\label{sec:risk_field_construction}
\Cref{alg:riskfield} summarizes the full two-model inference pipeline: the single-target flow gives the ego occupancy $P_{\mathrm{ego}}$ (and its MAP path defines the action proxy $\mathbf{a}_{1:K}$); the scene-level autoregressive joint gives each surrounding agent's ego-conditioned occupancy on the grid via the MAP rollout; agent velocities come from the predicted density's centroid motion; meeting probability is the oriented-box overlap; and the kinetic-energy severity weights it, accumulated over agents inside the outer counterfactual loop over candidate ego maneuvers. Both flows are conditioned on the map.

\begin{algorithm}[!t]
\caption{World-model-conditioned counterfactual spatial-temporal risk field}
\label{alg:riskfield}
\begin{algorithmic}[1]
\REQUIRE history $\mathbf{X}$; map $\mathbf{M}$; surrounding vehicles $\mathcal{J}$; grid $\mathcal{G}$ ($S{\times}S$); horizon $K$; candidate ego maneuvers $\{\mathcal{A}^{(1)},\dots,\mathcal{A}^{(L)}\}$; trained ego flow $\theta$ and scene-level joint $\psi$
\ENSURE risk tensors $\{\mathbf{R}(\mathcal{A}^{(\ell)})\}_{\ell=1}^{L}$; recommended $\mathcal{A}^{\star}$
\FORALL{candidate maneuver $\mathcal{A}=\mathbf{a}_{1:K}$}
  \STATE $P_{\mathrm{ego},1:K}\leftarrow$ single-target flow $\theta$ on $\mathcal{G}$ (map-conditioned)
  \STATE $\mathbf{v}^{(0)}\leftarrow$ centroid velocity of $P_{\mathrm{ego}}$; \ $\mathbf{a}_{1:K}\!\leftarrow$ proxy of $P_{\mathrm{ego}}$ MAP (or supplied $\mathcal{A}$)
  \STATE $\mathbf{s}^{\mathrm{sc}}_{1:K}\leftarrow$ scene rollout driven by $\mathbf{a}_{1:K}$; \ ordering $\pi$; \ $\mathbf{Y}_{\pi(0)}\!\leftarrow P_{\mathrm{ego}}$ MAP
  \STATE initialize $\mathrm{Risk}^{\mathrm{all}}_{k}(\cdot)\!\leftarrow\!0$
  \FORALL{agent $j\in\mathcal{J}$ in order $\pi$}
    \STATE $P^{(j)}_{1:K}\leftarrow$ joint conditional $p_\psi(\cdot\mid\mathbf{Y}_{\pi(<i)},\mathbf{s}^{\mathrm{sc}})$ on $\mathcal{G}$; \ fill $\mathbf{Y}_j\!\leftarrow$ its MAP
    \STATE $\mathbf{v}^{(j)}\leftarrow$ centroid velocity of $P^{(j)}$
    \FOR{$k=1$ \TO $K$}
      \STATE $M^{(j)}_k(\mathbf{g})\leftarrow$ oriented-box overlap of $P^{(j)}_k$ \COMMENT{\Cref{eq:box_overlap}}
      \STATE $C^{(j)}_{k}\leftarrow\tfrac12 m^{(j)}_{\mathrm{r}}\lVert\mathbf{v}^{(0)}_{k}-\mathbf{v}^{(j)}_{k}\rVert_2^2$ \COMMENT{per-step scalar}
      \STATE $\mathrm{Risk}^{\mathrm{all}}_{k}\mathrel{+}= C^{(j)}_{k}\,\big(P_{\mathrm{ego},k}\odot M^{(j)}_{k}\big)$
    \ENDFOR
  \ENDFOR
  \STATE $\mathbf{R}(\mathcal{A})\leftarrow\mathrm{stack}_{k=1}^{K}\,\mathrm{Risk}^{\mathrm{all}}_{k}$ \COMMENT{$S^2\times K$ tensor}
  \STATE $J(\mathcal{A})\leftarrow\sum_{k}\gamma^{k}\sum_{\mathbf{g}}\mathrm{Risk}^{\mathrm{all}}_{k}(\mathbf{g}\mid\mathcal{A})$
\ENDFOR
\STATE $\mathcal{A}^{\star}\leftarrow\arg\min_{\mathcal{A}}J(\mathcal{A})$
\RETURN $\{\mathbf{R}(\mathcal{A}^{(\ell)})\}$,\ $\mathcal{A}^{\star}$
\end{algorithmic}
\end{algorithm}

Equivalently, for a candidate ego maneuver $\mathcal{A}$, stacking the per-step accumulated maps over $K$ steps yields the risk tensor
\begin{equation}
\mathbf{R}(\mathcal{A})\in\mathbb{R}^{S^2\times K},\qquad \mathbf{R}[i,k]=\mathrm{Risk}^{\mathrm{all}}_{k}(\mathbf{g}_i\mid\mathcal{A}).
\label{eq:risk_tensor}
\end{equation}
This representation preserves temporal evolution coherently (shared latent rollout), supports adjustable spatial resolution, keeps a physically interpretable density$\times$severity decomposition, and is counterfactual in the ego action.

\FloatBarrier
\section{Results} \label{sec:results}
We report (i) a comparison against representative risk-quantification methods, (ii) a temporal-coherence ablation isolating the world-model contribution, (iii) a counterfactual demonstration, and (iv) qualitative density/risk visualizations. Trajectory-prediction accuracy is not our contribution; we therefore evaluate the \emph{risk output}, not predictor leaderboards.

\subsection{Experimental Setup}\label{sec:results_setup}
\subsubsection{Dataset}\label{sec:results_dataset}
Experiments use the InD (Intersection Drone) dataset \cite{bock2020ind}, which features highly interactive intersection scenarios with multi-modal futures. We use \emph{all} 33 InD recordings spanning the four physically distinct intersections (locations). Because the four sites have different coordinate frames, spatial inputs are normalized \emph{per location} (each recording's positions mapped to $[0,1]$ by its own location bounding box), so a single model is trained across all intersections without coordinate leakage. Each sample is a sliding window of $T{=}50$ history and $K{=}50$ future frames at a sampling step of $2$ (i.e., $\Delta t{=}0.08$\,s), with up to $8$ agents. After windowing, a per-recording $75\%/25\%$ split yields $213{,}504$ training and $71{,}173$ test windows ($284{,}677$ total). Cross-\emph{location} generalization is therefore built into a single model. The model architecture and training configuration are summarized in \Cref{tab:setup}.

\begin{table}[h]
\caption{Model and training setup (InD).}
\label{tab:setup}
\centering
\begin{tabularx}{\columnwidth}{>{\centering\arraybackslash}X >{\centering\arraybackslash}X}
\toprule
\textbf{Item} & \textbf{Value} \\
\midrule
History length $T$ (frames) & 50 \\
Prediction horizon $K$ (frames) & 50 \\
Batch size & 64 \\
Epochs & 50 \\
Learning rate & $10^{-3}$ \\
World-model state dim $d_s$ & 256 \\
Aux.\ dynamics weight $\lambda_{\mathrm{dyn}}$ & 5.0 \\
Number of parameters & 23{,}229{,}786 \\
Training time (all-33, 1$\times$A40) & $\approx$9.2\,h \\
\bottomrule
\end{tabularx}
\end{table}

\subsubsection{Conflict Labels and Metrics}\label{sec:results_metrics}
To evaluate a \emph{risk} representation rather than a predictor, we label scene-level conflict events. For the ego and each surrounding agent $j$ we form per-frame \emph{oriented} vehicle boxes inflated by a near-miss clearance $\delta$, and extend both ground-truth trajectories by $\tau_{\mathrm{TTC}}$ at constant velocity so that closing approaches just beyond the horizon are caught. The pair \emph{meets} at frame indices $(k_i,k_j)$ iff (a) the inflated boxes overlap, decided by the separating-axis test, \emph{and} (b) the relative speed $\lVert\mathbf{v}^{(0)}_{k_i}-\mathbf{v}^{(j)}_{k_j}\rVert>v_{\min}$ --- the speed gate excludes static car-following, which is ubiquitous in dense traffic. The post-encroachment time is $\mathrm{PET}_j=\min_{(k_i,k_j)\,\mathrm{meet}}|k_i-k_j|\,\Delta t$, and a scene is a \emph{conflict} iff $\min_j \mathrm{PET}_j<\tau_{\mathrm{PET}}$. The \emph{primary} thresholds $(\delta,v_{\min},\tau_{\mathrm{PET}},\tau_{\mathrm{TTC}})=(0.3\,\mathrm{m},\,2.5\,\mathrm{m/s},\,1.0\,\mathrm{s},\,1.5\,\mathrm{s})$ give $13.4\%$ positives on InD ($11.4\%$ on AD4CHE). For robustness we also evaluate against a \emph{strict} label ($\delta{=}0.2$\,m, $v_{\min}{=}3.0$\,m/s; $7.6\%$ positives) and a \emph{TTC-independent} label (minimum centre-distance below $3$\,m anywhere in the horizon, with no boxes, speed gate, or extrapolation; $8.6\%$ positives). Crucially the labels depend only on ground-truth geometry. 

We compare the controlled \emph{same-backbone} family --- \textbf{Ours} (expected energy), \textbf{Ours (prob.)} (peak meeting probability, no severity), and \textbf{PORA-style} occupancy-overlap on the \emph{same} predicted densities \cite{wang2025dynamic} --- against the handcrafted kinematic \textbf{DSF} safety field. We report three metrics. \textbf{AUROC}, the area under the receiver-operating-characteristic (true-positive-rate vs.\ false-positive-rate) curve, equals the probability that a randomly chosen conflict scene receives a higher risk score than a randomly chosen non-conflict scene; it is threshold-free ($0.5$ is chance, $1$ perfect). \textbf{AP} (average precision) is the area under the precision--recall curve, $\mathrm{AP}=\sum_n (R_n-R_{n-1})\,P_n$ with $P_n$ and $R_n$ the precision and recall at the $n$-th score threshold; by focusing on the positive (conflict) class it is the more informative metric under the $8$--$13\%$ class imbalance, with a no-skill value equal to the positive rate. \textbf{Early-warning lead time} is the mean time by which a method's risk first exceeds an operating threshold ahead of the labelled conflict onset, the threshold fixed at a $10\%$ false-alarm rate on non-conflict scenes. Every method is evaluated under all three label definitions (\Cref{tab:robust}).

\subsection{Risk-Quantification Comparison}\label{sec:results_compare}
We evaluate each method as a binary detector of surrogate conflicts on a $2{,}561$-scene stratified sample of the all-33 test split. The results are in \Cref{tab:robust}.
Three findings, reported plainly. 
(i) \emph{Our field leads on the metrics that matter under imbalance}: \emph{ours} (prob.) attains the best AP and the best early-warning lead on \emph{every} label definition --- AP being the more informative metric under the $8$--$13\%$ class imbalance. On AUROC, which rewards coarse global ranking and is forgiving under heavy imbalance, the handcrafted DSF field is marginally ahead on the box-based primary and strict labels (e.g.\ $0.719$ vs.\ $0.716$), but it trails ours on AP by a wide margin ($0.256$ vs.\ $0.324$) and yields no probabilistic field, severity, or counterfactual; on the geometry-only TTC-independent label our field leads even on AUROC ($0.805$ vs.\ $0.800$).
(ii) \emph{The controlled same-backbone comparison is decisive}: on the identical predicted densities, our exact-density oriented-box field beats the PORA-style occupancy-overlap on \emph{every} label and \emph{both} metrics (AUROC $+0.05$, AP $+0.10$ here), isolating the value of the exact density and oriented footprint over a discretized overlap. Beyond these numbers, our representation provides a spatially continuous probabilistic field, a physical severity, and a counterfactual query --- outputs a scalar SSM cannot produce. 

\emph{Our field exposes two complementary reads of the same predicted density --- one for each question a risk field must answer.} The \emph{probability} read (peak meeting probability, \emph{ours-prob}) answers \emph{whether} a conflict occurs and is the read that delivers the detection headline above (best AP and best early-warning lead). The \emph{energy} read (probability $\times$ kinetic-energy severity, \emph{ours}) answers \emph{how severe} the conflict would be. We further evaluate the energy read against the ground-truth collision energy on the \emph{same} predicted densities (Spearman $\rho{=}0.26$, $p{<}10^{-6}$, over the $343$ conflict scenes) where the probability read does not ($\rho{=}0.03$, n.s.). The two reads are thus not competitors but the two questions a risk field must answer: \emph{will} it happen, and \emph{how bad}. 

\begin{table*}[t]
\caption{Conflict-detection comparison and label-definition robustness on InD ($2{,}561$ scenes). AUROC, AP, and early-warning lead time (s, at $10\%$ false-alarm rate; scalar SSMs are single-shot, ``---'') under the primary ($13.4\%$), strict ($7.6\%$), and TTC-independent distance ($8.6\%$) conflict labels. Best per column in bold: ours (prob.) tops AP and early-warning lead under all three labels (and AUROC on the TTC-independent label); the handcrafted DSF leads AUROC on the box-based labels.}
\label{tab:robust}
\centering
\footnotesize
\setlength{\tabcolsep}{4.5pt}
\begin{tabular}{l ccc ccc ccc}
\toprule
& \multicolumn{3}{c}{\textbf{Primary} ($13.4\%$)} & \multicolumn{3}{c}{\textbf{Strict} ($7.6\%$)} & \multicolumn{3}{c}{\textbf{TTC-indep.} ($8.6\%$)} \\
\cmidrule(lr){2-4}\cmidrule(lr){5-7}\cmidrule(lr){8-10}
\textbf{Method} & AUROC & AP & Lead & AUROC & AP & Lead & AUROC & AP & Lead \\
\midrule
DSF \cite{4,6}          & $\mathbf{0.719}$ & $0.256$ & --- & $\mathbf{0.711}$ & $0.140$ & --- & $0.800$ & $0.260$ & --- \\
PORA-style \cite{wang2025dynamic}  & $0.662$ & $0.221$ & $-0.02$ & $0.664$ & $0.158$ & $0.00$ & $0.720$ & $0.167$ & $0.32$ \\
\midrule
\textbf{Ours} (prob.)   & $0.716$ & $\mathbf{0.324}$ & $\mathbf{0.37}$ & $0.695$ & $\mathbf{0.241}$ & $\mathbf{0.51}$ & $\mathbf{0.805}$ & $\mathbf{0.325}$ & $\mathbf{0.68}$ \\
\bottomrule
\end{tabular}
\end{table*}

\subsection{Qualitative analysis}
\label{sec:results_qualitative}
\Cref{fig:qualitative} shows the field on two genuine conflicts of contrasting kinematics, each read through four lenses of the same field and overlaid on the recorded drone map. For every scene the panels show, left to right: the predicted per-agent occupancy densities $P^{(i)}(g,k^\star)$ at the highest-risk frame $k^\star$, shown \emph{ahead} of each vehicle's present position (oriented boxes mark the current positions; ego $\circ$ in cyan, all other agents $\square$ in orange; ground-truth tracks overlaid); each agent's centroid (bulk) velocity that feeds the severity (one arrow per agent, coloured by speed); the resulting risk field $R_k(g)$ at $k^\star$ as a heat overlay (expected collision energy); and the spatial peak $\max_g R_k(g)$ traced across the prediction horizon. \emph{Scene~A} (top) is a \emph{fast but near-parallel} approach: the ego and its neighbour move at $13$ and $11$\,m/s only $7^\circ$ apart, so their \emph{relative} speed is just $2.8$\,m/s and the risk peaks at a modest $22$\,J (PET $0.32$\,s). \emph{Scene~B} (bottom) is a \emph{slower but near-head-on} approach: $5$ and $7$\,m/s at $176^\circ$, giving a $12.1$\,m/s relative speed and an order-of-magnitude higher $671$\,J peak (PET $0.24$\,s). In both, the predicted densities overlap \emph{ahead} of the present positions and the heat overlay localizes on that forecast meeting point, well above the $10.0$\,J per-cell critical level. The contrast makes the severity physics explicit: the kinetic-energy term tracks the \emph{relative} velocity set by the encounter geometry, not absolute speed.

\begin{figure*}[t]
\centering
\subfigure[Scene A:predicted occupancy]{\includegraphics[height=2.9cm]{figures/qual_map_crit_pos.png}}\hfill
\subfigure[Scene A:velocity field]{\includegraphics[height=2.9cm]{figures/qual_map_crit_vel.png}}\hfill
\subfigure[Scene A:risk field]{\includegraphics[height=2.9cm]{figures/qual_map_crit_risk.png}}\hfill
\subfigure[Scene A:risk-field energy]{\includegraphics[height=2.9cm]{figures/qual_map_crit_energy.png}}\\[2pt]
\subfigure[Scene B:predicted occupancy]{\includegraphics[height=2.9cm]{figures/qual_map_ncrit_pos.png}}\hfill
\subfigure[Scene B:velocity field]{\includegraphics[height=2.9cm]{figures/qual_map_ncrit_vel.png}}\hfill
\subfigure[Scene B:risk field]{\includegraphics[height=2.9cm]{figures/qual_map_ncrit_risk.png}}\hfill
\subfigure[Scene B:risk-field energy]{\includegraphics[height=2.9cm]{figures/qual_map_ncrit_energy.png}}
\caption{Two InD conflicts of contrasting kinematics, overlaid on the recorded drone map. \textbf{Top (Scene A):} a fast but near-parallel approach ($13$/$11$\,m/s, $7^\circ$ apart; relative speed $2.8$\,m/s; peak $22$\,J). \textbf{Bottom (Scene B):} a slower but near-head-on approach ($5$/$7$\,m/s, $176^\circ$; relative speed $12.1$\,m/s; peak $671$\,J). Each row shows, left to right: the predicted per-agent occupancy densities at the highest-risk frame, shown ahead of each vehicle (oriented boxes mark the present positions; ego $\circ$ cyan, others $\square$ orange; ground-truth tracks overlaid), each agent's centroid (bulk) velocity underlying the severity (one arrow per agent, coloured by speed), the resulting risk field at that frame as a heat overlay (expected collision energy), and the spatial peak of the risk field over the horizon (per-scene axes). The severity tracks \emph{relative} kinetic energy: high for the head-on pair despite its lower absolute speed.}
\label{fig:qualitative}
\end{figure*}

\section{Ablation and Analysis}\label{sec:results_analysis}
This section isolates the contributions of the world model and map conditioning, and reports a counterfactual demonstration. We also report qualitative visualizations in \Cref{sec:results_qualitative}.

\subsection{Ablation Across Model Variants}\label{sec:results_ablation}
We compare four variants (\Cref{tab:ablation}) on the full all-33 split under an identical pipeline, training schedule, and budget, toggling only the conditioning and training scheme. \emph{No WM} conditions the flow on the broadcast context with the flow's internal time index (the $[\mathbf{c},k]$ marginal model). \emph{Naive WM} replaces it with an unconstrained latent rollout. \emph{WM, residual} adds residual conditioning and the auxiliary dynamics loss. \emph{Multi-agent action-conditioned (ours)} additionally trains the model to predict a surrounding agent's future conditioned on the ego's realized action proxy. Accuracy uses NLL/CRPS/minADE/minFDE/RMSE on the $71{,}173$-window test set ($100$ samples). CRPS (continuous ranked probability score) is the integral over the support of the squared difference between the predicted and the empirical (Heaviside) cumulative distributions; it jointly rewards calibration and sharpness and reduces to the absolute error for a deterministic prediction (lower is better). Temporal coherence is measured on mode-consistent trajectories (one base code held fixed and decoded through the horizon under each step's conditioning) by the \emph{Acc.\ ratio}: mean trajectory acceleration relative to ground truth (closer to $1$ is better).

\begin{table}[h]
\caption{Ablation across model variants on the all-33 InD test split. Accuracy: lower is better. Coherence: Acc.\ ratio closer to $1$ is better. Best in each column in bold.}
\label{tab:ablation}
\centering
\footnotesize
\setlength{\tabcolsep}{3pt}
\begin{tabular}{l c c c c}
\toprule
\textbf{Variant} & \textbf{NLL} & \textbf{CRPS} & \textbf{mADE/mFDE/RMSE} & \textbf{Acc.} \\
\midrule
No WM ($[\mathbf{c},k]$)        & $-6.745$ & $0.730$ & $0.248/0.824/1.789$ & $5.40$ \\
Naive WM                        & $-3.583$ & $3.408$ & $1.043/2.083/8.700$ & $5.94$ \\
WM, residual                    & $-6.793$ & $0.733$ & $0.241/0.774/1.851$ & $5.64$ \\
\textbf{Multi-agent (ours)}     & $\mathbf{-6.948}$ & $\mathbf{0.694}$ & $\mathbf{0.236/0.834/1.655}$ & $\mathbf{4.54}$ \\
\bottomrule
\end{tabular}
\end{table}

Three findings, reported plainly. (i) \emph{A naive world model hurts}: an unconstrained latent rollout collapses accuracy ($\mathrm{NLL}\,{-}3.58$ vs.\ ${-}6.74$, $\mathrm{minADE}\,1.04$ vs.\ $0.25$) --- the rollout becomes a weak conditioning bottleneck and underfits; residual conditioning plus the auxiliary dynamics loss fixes this. (ii) \emph{Multi-agent action-conditioned training is the best variant}: predicting a surrounding agent conditioned on the ego's action proxy attains the best NLL, CRPS, minADE, and RMSE, and the best Acc.\ ratio ($4.54$, closest to real-driving smoothness), of all variants. The richer multi-agent objective acts as a beneficial auxiliary signal rather than a cost --- only minFDE is marginally behind the residual-WM variant. (iii) \emph{Coherence improvements are real but modest}: trajectory roughness (Acc.\ ratio) improves --- ours is closest to ground truth --- yet all variants remain several times rougher than real driving in acceleration, so we do not over-claim a large coherence gain.

\subsection{Physically Grounded Counterfactual Risk}\label{sec:results_counterfactual}
We evaluate the counterfactual objective $J(\mathcal{A})$ (\Cref{eq:counterfactual_objective}) under candidate ego maneuvers (maintain, brake, accelerate). The counterfactual signal is the \emph{physically grounded} ego severity term: a candidate maneuver determines the ego velocity profile, which changes the kinetic-energy-loss cost $C^{(j)}_k$ (\Cref{eq:collision_cost}) and hence the aggregated risk. This counterfactual is causal by construction --- braking genuinely lowers collision kinetic energy by the laws of mechanics. In the two-model field the surrounding density is additionally ego-conditioned through the joint, so a learned channel is also present; however, as \Cref{sec:results_negative} shows, on observational data it is weak, so the realized counterfactual remains dominated by the physical severity term. \Cref{tab:cf} reports $J(\mathcal{A})$ on the two-model field, averaged over $20$ test scenes ($\gamma=0.95$), where each maneuver scales the ego speed profile by $\pm30\%$; the safest maneuver $\mathcal{A}^\star$ is the argmin.

\begin{table}[h]
\caption{Physically grounded counterfactual objective $J(\mathcal{A})$ (\Cref{eq:counterfactual_objective}) per ego maneuver on the two-model field, mean over $20$ test scenes ($\gamma=0.95$, $\pm30\%$ ego-speed maneuver). Lower is safer. The response is quadratic in relative speed, hence asymmetric.}
\label{tab:cf}
\centering
\footnotesize
\begin{tabularx}{\columnwidth}{l c c}
\toprule
\textbf{Ego maneuver} $\mathcal{A}$ & $J(\mathcal{A})$ & \textbf{vs.\ maintain} \\
\midrule
Maintain speed          & $2.63\times10^{4}$ & --- \\
\textbf{Brake} ($\mathcal{A}^\star$) & $\mathbf{1.88\times10^{4}}$ & $-28.3\%$ \\
Accelerate              & $3.69\times10^{4}$ & $+40.4\%$ \\
\bottomrule
\end{tabularx}
\end{table}

The ranking is interpretable and physically correct: braking lowers the relative kinetic energy and is selected as $\mathcal{A}^\star$, while accelerating raises it. Fixed-plan occupancy methods such as PORA \cite{wang2025dynamic} score only a single supplied plan and cannot produce this comparison. This is the axis on which the kinetic-energy severity earns its place: it is deliberately a \emph{harm} weight rather than an occurrence cue --- hence it trails the pure-probability read on conflict detection (\Cref{sec:results_compare}) --- and it is precisely what makes the maneuver ranking physically meaningful.

\subsection{A Learned Reactive Counterfactual Does Not Emerge: A Pre-Registered Negative Result}\label{sec:results_negative}
The counterfactual of \Cref{sec:results_counterfactual} reshapes the \emph{severity} term. A stronger capability would be a \emph{learned reactive} counterfactual, in which the predicted \emph{occupancy density} of surrounding agents itself responds to the ego's action (others react to the ego's plan). We tested this with pre-registered criteria across three escalating designs: (a) conditioning the ego prediction on a neighbour's action; (b) the multi-agent action-conditioned model (\Cref{sec:results_ablation}); and (c) an architecturally correct scene-level \emph{autoregressive joint} model that predicts all agents jointly, $p(\mathbf{Y}_{1:N}\mid\text{scene},\mathbf{a}_{\mathrm{ego}})=\prod_i p(\mathbf{Y}_i\mid\mathbf{Y}_{<i},\text{scene},\mathbf{a}_{\mathrm{ego}})$, with the ego explicitly identified.

We decompose $J(\mathcal{A})$ into $J_{\mathrm{full}}$ (density and cost both react), $J_{\mathrm{dens}}$ (only the learned density reacts), and $J_{\mathrm{phys}}$ (only the physical cost reacts). The pre-registered success criteria were: per-agent density shift $\mathrm{L1}>0.01$; $J_{\mathrm{dens}}$ spread $>5\%$ across maneuvers; and $J_{\mathrm{full}}$/$J_{\mathrm{phys}}$ rankings differing on $\geq30\%$ of scenes. \Cref{tab:negative} reports the scene-level model.

\begin{table}[h]
\caption{Learned-counterfactual diagnostic for the scene-level autoregressive model. All three pre-registered criteria fail.}
\label{tab:negative}
\centering
\footnotesize
\begin{tabularx}{\columnwidth}{l c c}
\toprule
\textbf{Quantity} & \textbf{Observed} & \textbf{Criterion} \\
\midrule
Density shift L1 (brake / accel) & $0.0046$ / $0.0047$ & $>0.01$ \quad\ding{55} \\
$J_{\mathrm{dens}}$ spread across maneuvers & $0.04\%$ & $>5\%$ \quad\ding{55} \\
$J_{\mathrm{full}}$ vs.\ $J_{\mathrm{phys}}$ agreement & identical to $0.01\%$ & differ \quad\ding{55} \\
\bottomrule
\end{tabularx}
\end{table}

To separate \emph{model quality} from \emph{counterfactual fidelity}, we add two direct tests on the map-conditioned models that back the risk field (\Cref{tab:wmfidelity}). First, the world model is \emph{factually accurate}: conditioned on the \emph{realized} ego action it attains a held-out neighbour conditional NLL of $-5.40$ and an ego marginal NLL of $-7.62$, so the negative result is not an artifact of an underfit model. Second, a counterfactual-fidelity test asks directly whether the ego action is \emph{predictively informative} about the neighbours: on scenes where the ego genuinely maneuvered (top-quartile realized $|\mathbf{a}_{\mathrm{ego}}|$), conditioning the joint on the \emph{true} ego action predicts neighbours' realized futures no better than conditioning on the \emph{reversed} action or a \emph{zero} (autonomous) action --- the neighbour-NLL gap is $+0.0004$ nats with a $95\%$ bootstrap CI of $[-0.0010,+0.0018]$, statistically indistinguishable from zero. Knowing what the ego actually did carries essentially no information about how the others responded: the conditional is well-estimated yet non-causal.

\begin{table}[h]
\caption{World-model factual quality vs.\ counterfactual fidelity. The model predicts the realized future well (top), yet the ego action is not predictively informative about neighbours' realized futures (bottom): the true-vs-counterfactual neighbour-NLL gap is indistinguishable from zero.}
\label{tab:wmfidelity}
\centering
\footnotesize
\begin{tabularx}{\columnwidth}{l c}
\toprule
\textbf{Quantity (held-out)} & \textbf{Value} \\
\midrule
Ego marginal NLL (realized action)        & $-7.62$ \\
Neighbour conditional NLL (realized action) & $-5.40$ \\
\midrule
$\Delta$NLL(reversed $-$ true), maneuver scenes & $+0.0004$ \ {\footnotesize$[-0.0010,+0.0018]$} \\
$\Delta$NLL(zero $-$ true), maneuver scenes     & $+0.0002$ \ {\footnotesize$[-0.0005,+0.0008]$} \\
\bottomrule
\end{tabularx}
\end{table}

All three pre-registered criteria fail: $J_{\mathrm{dens}}$ is flat across maneuvers and $J_{\mathrm{full}}\!\approx\!J_{\mathrm{phys}}$ to within $0.01\%$ --- the counterfactual signal is essentially $100\%$ physics. The learned occupancy density is non-reactive to the ego action. The map-conditioned scene-level model used for the risk field corroborates this: under a \emph{realistic} maneuver (perturbation scaled to the agent's own action magnitude) the per-agent density shift is only $\mathrm{L1}\!\approx\!0.002$, still below the $0.01$ bar, and the response rises smoothly with action magnitude --- reaching the large shifts ($\mathrm{L1}\!\approx\!0.48$) only at \emph{out-of-distribution} ego actions. The mechanism is therefore present and correctly wired (unlike earlier designs where it was flat at zero), but the realistic-maneuver effect is bounded by the data, not the architecture. We attribute this to the \emph{observational/interventional gap}: InD records what happened, never what would have happened under a different ego action, so training on observed $(\mathbf{a}_{\mathrm{ego}},\mathbf{Y})$ pairs learns a conditional that conflates correlation with confounding, not the interventional response. No architecture overcomes this from observational logs alone. The physical counterfactual of \Cref{sec:results_counterfactual} works precisely because mechanics is a \emph{causal} model that is encoded, not learned. A learned reactive counterfactual would require interventional or simulation data (\Cref{section:conclusion_future_work}).

\subsection{Cross-Dataset Generalization: AD4CHE Congested Highway}\label{sec:results_ad4che}
To test that the method is not InD-specific, we retrain and evaluate the \emph{entire} two-model field on AD4CHE~\cite{zhang2023ad4che}, an aerial dataset of \emph{congested} expressway traffic --- structurally very different from InD's intersections (dense car-following at highway speed vs.\ crossing/turning at low speed). We use the dataset's 11 recording sites with \emph{per-scene} spatial normalization (analogous to InD's per-location scheme) and the per-scene road raster for map conditioning; the only dataset-specific code is the loader. Windowing matches InD ($T{=}K{=}50$, ${\sim}3.3$\,s horizon at 30\,fps), giving $11{,}707$ training and $3{,}934$ test windows.

\paragraph{Prediction quality.} The models train cleanly on AD4CHE: held-out ego-marginal NLL $-6.43$ and neighbour-conditional NLL $-5.11$, comparable to InD ($-7.57$, $-5.44$); congested highway is denser and slightly less sharp.

\paragraph{Conflict detection.} \Cref{tab:ad4che_detect} repeats the comparison of \Cref{tab:robust} on AD4CHE conflict labels ($11.4\%$ positives). Here our field wins \emph{decisively on AUROC, AP, and early-warning lead}: ours (prob.) attains AUROC $\mathbf{0.705}$ / AP $\mathbf{0.357}$ and a $0.31$\,s lead, well clear of the handcrafted DSF field ($0.584$ / $0.148$) and PORA-style (whose lead is negative, $-0.49$\,s). The contrast with InD is informative: on InD the handcrafted DSF field was competitive on AUROC, but in stop-and-go congestion constant-velocity assumptions break down (vehicles continually accelerate/brake), so the kinematic baseline collapses while the learned world-model field captures the dynamics. The controlled same-backbone advantage (ours $>$ PORA-style) holds on both metrics here too.

\begin{table}[h]
\caption{Conflict detection on AD4CHE ($394$ scenes, $11.4\%$ positives). AUROC, AP, and early-warning lead time (s, at $10\%$ false-alarm rate; single-shot SSMs ``---''). Higher is better; best per column in bold.}
\label{tab:ad4che_detect}
\centering
\footnotesize
\setlength{\tabcolsep}{4pt}
\begin{tabular}{l c c c}
\toprule
\textbf{Method} & \textbf{AUROC} & \textbf{AP} & \textbf{Lead} \\
\midrule
DSF \cite{4,6}               & $0.584$ & $0.148$ & --- \\
PORA-style \cite{wang2025dynamic}       & $0.651$ & $0.209$ & $-0.49$ \\
\midrule
\textbf{Ours} (prob.)        & $\mathbf{0.705}$ & $\mathbf{0.357}$ & $\mathbf{0.31}$ \\
\bottomrule
\end{tabular}
\end{table}

\paragraph{Per-dataset calibration.} The per-cell critical level $T_{\mathrm{critical}}$ differs by ${\sim}18\times$ across datasets ($10.0$\,J on InD vs.\ $183$\,J on AD4CHE), because highway speeds produce far larger collision energies; each dataset is therefore calibrated independently. \Cref{fig:ad4che} shows the AD4CHE field on its own scale --- risk concentrating where an ego drives onto a queue of near-static vehicles in dense highway traffic.

\paragraph{The negative result generalizes.} The learned-counterfactual proxy (\Cref{sec:results_negative}) gives the same null on AD4CHE (true-vs-counterfactual neighbour-NLL gap ${\approx}0$, CI spanning zero), confirming the observational/interventional gap is dataset-general rather than an InD artifact.

\begin{figure}[t]
\centering
\subfigure[Site 17: ego truck (E) closing on a stopped truck queue]{\includegraphics[width=0.92\columnwidth]{figures/qual_map_c530_risk.png}}\\[2pt]
\subfigure[Site 15: ego car (E) closing on a stopped truck]{\includegraphics[width=0.92\columnwidth]{figures/qual_map_c190_risk.png}}
\caption{Two AD4CHE conflicts on the dataset's own calibration ($T_{\mathrm{critical}}{=}183$\,J colorbar). Oriented boxes mark each vehicle at the initial frame using its recorded dimensions and heading; labels give the agent id and observed speed in m/s (ego ``E'' in cyan, other agents in orange). The heat overlay is the risk field at its peak-risk frame over the horizon, as in \Cref{fig:qualitative}. In both, the ego's predicted occupancy drives onto a cluster of near-static vehicles (high meeting probability) while its own relative speed supplies the severity, so the field rises above critical. \Cref{tab:ad4che_decomp} decomposes the peak cell of each.}
\label{fig:ad4che}
\end{figure}

\begin{table}[t]
\caption{Why each ego in \Cref{fig:ad4che} is high-risk: decomposition of the risk field at its peak cell $\mathbf{g}^\star$, $\mathrm{Risk}(\mathbf{g}^\star)=\underbrace{P_{\mathrm{ego}}\cdot(P_j\!\circledast\!\mathbf{1}_{\mathrm{box}})}_{\text{meeting prob.\ }M}\cdot\underbrace{\tfrac12\mu\|\Delta\mathbf{v}\|^2}_{\text{severity }C}$, with $\Delta\mathbf{v}$ the bulk (centroid) relative velocity, $\mu{=}750$\,kg, and $T_{\mathrm{critical}}{=}183$\,J. The dominant conflict partner $j$ is shown; both peaks exceed critical.}
\label{tab:ad4che_decomp}
\centering
\setlength{\tabcolsep}{5pt}
\begin{tabular}{l c c}
\toprule
 & \textbf{Site 17 (a)} & \textbf{Site 15 (b)} \\
\midrule
Ego (class, speed m/s)            & truck, $4.1$       & car, $9.1$ \\
Partner $j$ (class, speed m/s)    & truck, $3.9$       & truck, $0.1$ \\
$P_{\mathrm{ego}}(\mathbf{g}^\star)$               & $0.10$             & $0.12$ \\
overlap $(P_j\!\circledast\!\mathbf{1}_{\mathrm{box}})(\mathbf{g}^\star)$ & $0.99$ & $0.99$ \\
meeting prob.\ $M(\mathbf{g}^\star)$               & $0.10$             & $0.12$ \\
rel.\ speed $\|\Delta\mathbf{v}\|$ (m/s)           & $6.1$              & $10.9$ \\
severity $C=\tfrac12\mu\|\Delta\mathbf{v}\|^2$ (J) & $1.4{\times}10^{4}$ & $4.4{\times}10^{4}$ \\
\textbf{peak risk} $M\!\cdot\!C$ (J)               & $\mathbf{1.4{\times}10^{3}}$ & $\mathbf{5.2{\times}10^{3}}$ \\
\bottomrule
\end{tabular}
\end{table}





\section{Conclusion and Future Work}\label{section:conclusion_future_work}
We presented a world-model-conditioned normalizing-flow risk field for safety assessment in multi-agent traffic. The field is ego-centric, $\mathrm{Risk}_{k}(\mathbf{g})=P_{\mathrm{ego},k}(\mathbf{g})\,M^{(j)}_{k}(\mathbf{g})\,C^{(j)}_{k}(\mathbf{g})$, and our method (i) uses an exact, grid-evaluable flow density instead of a discretized heatmap; (ii) conditions that density on a residual deterministic latent world-model rollout and on a bird's-eye map of the road geometry, which we show is load-bearing for accuracy; (iii) couples two such models --- a single-target flow for the ego's occupancy and a scene-level autoregressive joint for the surrounding agents' occupancies conditioned on the ego --- with meeting probability from oriented-box overlap and severity from a kinetic-energy-loss term whose relative velocity is read from the predicted density's centroid motion; and (iv) supports a physically grounded counterfactual risk query --- ranking candidate ego maneuvers by collision risk --- that fixed-plan occupancy methods such as PORA \cite{wang2025dynamic} structurally lack.

Equally, we report a careful negative result. Using pre-registered criteria, we found that a \emph{learned reactive} counterfactual --- surrounding agents' predicted occupancy responding to the ego's action --- does not emerge from observational driving logs, even with an architecturally correct scene-level autoregressive joint model. This is the observational/interventional gap: an observational dataset records what happened, not what would have happened under a different ego action, so it cannot identify the interventional response. The physically grounded counterfactual works precisely because mechanics is a causal model that is encoded rather than learned.

The clearest path to a \emph{learned} reactive counterfactual is therefore interventional data. Future work will train a world model on large-scale and closed-loop sources --- e.g.\ nuPlan and microsimulation --- where the ego action can be intervened upon and the reaction observed, which is the prerequisite the present observational study identifies. Additional directions include a stochastic latent world model for explicit mode branching, executing the risk-quantification baseline comparison protocol, closed-loop evaluation with a risk-aware planner, and cross-dataset generalization (highD/rounD).

\bibliographystyle{IEEEtran}
\bibliography{reference}
\end{document}
